% ===========================================================================
\chapter{Talking to Claude}\label{ch:claude}

Everything so far has worked with a pretend model. Now the real one. The
agent talks to Claude through Anthropic's \emph{Messages API}: one HTTP
request per reply, to \url{https://api.anthropic.com/v1/messages}, with
the conversation, the tools and the system prompt as JSON. The answer is a
JSON reply, of the shape we have been writing by hand since
Chapter~\ref{ch:conversation}. So \code{Claude_Model} is mostly about
HTTP: building the request, sending it, and coping when it fails.

\section{The request}

Here is a request with everything we send, abridged:

\begin{json}
POST /v1/messages
content-type: application/json
x-api-key: sk-ant-...
anthropic-version: 2026-10-21
anthropic-beta: server-side-fallback-2026-07-01

{"model": "claude-opus-5-5",
 "max_tokens": 16000,
 "system": "You are a coding agent working in a project folder ...",
 "tools": [{"name": "list_files", ...}, {"name": "read_file", ...}, ...],
 "messages": [{"role": "user", "content": [...]}, ...],
 "cache_control": {"type": "ephemeral"},
 "output_config": {"effort": "high"},
 "fallbacks": "default"}
\end{json}

The headers say who we are (the API key) and which version of the API we
are written for. In the body:

\begin{itemize}
\item \code{model} is which Claude. We use \code{claude-opus-5-5}; any
  other can be chosen on the command line.
\item \code{max_tokens} is the longest reply allowed. A reply that reaches
  it is cut off, with the \code{stop_reason} \code{"max_tokens"}, which the
  agent handles (Chapter~\ref{ch:loop}).
\item \code{system}, \code{tools} and \code{messages} are what a
  \code{Request} holds.
\item \code{cache_control} asks the API to cache the request. Every
  request repeats the one before it, plus a little; with caching, that
  repeated part is read from the cache at about a tenth of the price, and
  faster. For an agent, which sends dozens of requests per task, each
  longer than the last, this is the difference between cheap and
  expensive.
\item \code{output_config.effort} says how hard the model should think:
  from \code{low} to \code{max}. Claude thinks before it answers, and more
  effort means more thinking, better work on hard problems, and more
  tokens. \code{high} is a good default for coding.
\item \code{fallbacks} deals with something rare. Claude has safety checks
  that may decline a request, now and then by mistake (code about
  security, say). With \code{"fallbacks": "default"}, and the
  \code{anthropic-beta} header that goes with it, the API then tries the
  request again on another model it recommends, instead of stopping. Not
  every model has fallbacks, so we only ask for them on those that do.
\end{itemize}

\section{The model}

A \code{Claude_Model} holds a \emph{credential}: either an API key, or a
login token (from Anthropic's \code{ant auth login}), which goes in a
different header. The usual way to make one is from the environment, so
that the key is never written in a file:

\begin{newfile}{claude_model.nex}
-- claude_model.nex: asking Claude, through the Messages API.
--
-- One HTTP request per reply: POST /v1/messages with the whole
-- conversation, the tools and the instructions as JSON. The answer is
-- JSON too, and becomes a Reply.

intern Model
intern Helpers
intern net/Http_Client

class Claude_Model inherit Model
create
  -- A model that signs its requests with an API key, or with a login
  -- token if `is_token`.
  make(credential: String, is_token: Boolean, model_id, effort: String)
  require
    has_credential: credential /= ""
    known_effort: Claude_Model.EFFORTS.contains(effort)
  do
    this.credential := credential
    this.is_token := is_token
    this.model_id := model_id
    this.effort := effort
    base_url := "https://api.anthropic.com"
  end

  -- A model using the credentials in the environment: ANTHROPIC_API_KEY,
  -- or else ANTHROPIC_AUTH_TOKEN. ANTHROPIC_BASE_URL, if set, sends the
  -- requests somewhere else (a proxy, or a test server).
  from_environment(model_id, effort: String)
  require
    has_credential: has_credentials()
    known_effort: Claude_Model.EFFORTS.contains(effort)
  do
    if program.has_env("ANTHROPIC_API_KEY") and program.env("ANTHROPIC_API_KEY") /= "" then
      make(program.env("ANTHROPIC_API_KEY"), false, model_id, effort)
    else
      make(program.env("ANTHROPIC_AUTH_TOKEN"), true, model_id, effort)
    end
    if program.has_env("ANTHROPIC_BASE_URL") and program.env("ANTHROPIC_BASE_URL") /= "" then
      base_url := program.env("ANTHROPIC_BASE_URL")
    end
  end

feature
  API_VERSION = "2026-10-21"
  MAX_TOKENS = 16000       -- the longest reply, in tokens
  EFFORTS = ["low", "medium", "high", "xhigh", "max"]
  ATTEMPTS = 5             -- tries per reply, when the API is busy
  TIMEOUT_MS = 600000      -- ten minutes: a hard turn can take a while

  credential: String
  is_token: Boolean
  model_id: String     -- which Claude: "claude-opus-5-5"
  effort: String       -- how hard it thinks: "low" .. "max"
  base_url: String

  name(): String do
    result := model_id
  end

  -- Send requests to `url` instead of the Claude API: a proxy, or a
  -- test server.
  set_base_url(url: String) do
    base_url := url
  end
\end{newfile}

\code{from_environment} calls \code{make} to do most of its work: one
constructor may call another, as an ordinary routine.
\code{Claude_Model.EFFORTS} reads a constant through the class's name,
which works even before there is an object.

\section{Sending, and trying again}

Requests over a network fail. Most failures are temporary: the API is
overloaded (status 529), we are sending requests too fast for our account
(429), the server had a problem (500), or the connection dropped. The
right response to all of these is to wait a little and try again, waiting
longer each time (2, 4, 8 seconds, \ldots), so as not to add to the
problem. If the API says how long to wait, in a \code{retry-after} header,
we wait that long.

Other failures are not temporary. A malformed request (400) or a wrong key
(401) will fail the same way every time, and trying again only wastes time.
Those stop at once, with the API's own explanation.

\begin{continued}{claude_model.nex}
  reply(request: Request): Reply do
    let client: Http_Client := create Http_Client.make()
    let body: String := request_body(request)
    let attempt: Integer := 0
    let answer: ?Reply := nil
    from until answer /= nil do
      attempt := attempt + 1
      let wait: Integer := retry_delay_for(attempt)
      let trouble: String := ""
      do
        let response: Http_Response := client.request("POST", base_url + "/v1/messages",
          body, headers(), Claude_Model.TIMEOUT_MS)
        if response.status() = 200 then
          answer := create Reply.from_json(json_value(response.body()))
        elseif is_retryable(response.status()) and attempt < Claude_Model.ATTEMPTS then
          trouble := "the API answered " + response.status().to_string
          wait := retry_delay(response, attempt)
        else
          raise api_error(response)
        end
      rescue
        -- A network failure (no connection, a timeout) is worth another
        -- try too; an error from the API itself is not.
        let message: String := "" + exception
        if message.starts_with("Claude API") or attempt >= Claude_Model.ATTEMPTS then
          raise message
        end
        trouble := preview(message, 80)
      end
      if answer = nil then
        console.print_line("  (" + trouble + "; trying again in "
          + count_text(wait / 1000, "second") + ")")
        sleep(wait)
      end
    end
    if ?answer as got then
      result := got
    else
      raise "no reply"    -- not reached: the loop ends only with one
    end
  end
\end{continued}

The structure takes a moment to read. Each time round the loop sends the
request once. A good answer ends the loop. A temporary failure sets
\code{trouble}, and the loop waits and goes round again. Anything that
raises an exception inside the \code{do} part (the network failing, or our
own \code{raise} for an error that is not temporary) lands in the
\code{rescue} part, which decides: an API error, or the last attempt, is
raised again and ends the routine; a network failure is one more reason
to wait and try again.

\code{client.request} is \code{Http_Client}'s way of sending a request
with headers of our choosing. The last argument is a timeout: a hard
problem, at high effort, can take the model minutes to answer, so we give
it ten.

The \code{if ?answer as got} at the end is for Nex's benefit. We know the
loop only ends once \code{answer} is set, but Nex can't see that, and
won't let us return a \code{?Reply} as a \code{Reply} without checking.

\section{The body and the headers}

The body is built as a Nex map, as in the request above, and turned into
JSON. The headers are a map too.

\begin{continued}{claude_model.nex}
  -- The JSON body of the request for `request`.
  request_body(request: Request): String do
    let body: Map[String, Any] := {
      "model": model_id,
      "max_tokens": Claude_Model.MAX_TOKENS,
      "system": request.system,
      "tools": request.tools,
      "messages": request.conversation.to_json(),
      -- Cache the conversation: the next request repeats all of it, and
      -- reading it from the cache costs a tenth as much.
      "cache_control": {"type": "ephemeral"}
    }
    let output_config: Map[String, Any] := {"effort": effort}
    body.put("output_config", output_config)
    if uses_fallbacks() then
      -- If a safety check declines the request, let the API try it
      -- again on the model it recommends, instead of stopping.
      body.put("fallbacks", "default")
    end
    result := json_text(body)
  end

  -- The HTTP headers every request carries.
  headers(): Map[String, String] do
    result := {"content-type": "application/json",
      "anthropic-version": Claude_Model.API_VERSION}
    let betas: Array[String] := []
    if is_token then
      result.put("authorization", "Bearer " + credential)
      betas.add("oauth-2025-04-20")
    else
      result.put("x-api-key", credential)
    end
    if uses_fallbacks() then
      betas.add("server-side-fallback-2026-07-01")
    end
    if not betas.is_empty then
      result.put("anthropic-beta", ",".join(betas))
    end
  end

  -- Whether the model can fall back to another when it declines.
  uses_fallbacks(): Boolean do
    result := ["claude-opus-5-5", "claude-opus-5", "claude-fable-5-1", "claude-fable-5",
      "claude-sonnet-5-5"].contains(model_id)
  end
\end{continued}

\code{output_config} is built on its own line, with its type written out,
before it goes into \code{body}. A map written straight into a call's
argument, as \code{body.put("output_config", \{"effort": effort\})}, would
be a \code{Map[String, String]}, and Nex does not treat that as a
\code{Map[String, Any]}: a map you can put anything into is not the same as
one that only holds strings.

\section{Errors}

Last, the rules for errors: which are worth another try, how long to wait,
and how to turn an error from the API into a sentence. The API's errors
are JSON, \code{\{"type": "error", "error": \{"type": ..., "message":
...\}\}}, and its message is usually clear enough to act on.

\begin{continued}{claude_model.nex}
  -- Whether an error is worth trying again after a pause: too many
  -- requests (429), the API overloaded (529) or failing (5xx), a timeout
  -- (408) or a conflict (409). Anything else, such as a malformed
  -- request (400) or a bad key (401), would fail again.
  is_retryable(status: Integer): Boolean do
    result := status = 408 or status = 409 or status = 429 or status >= 500
  end

  -- How long to wait before attempt `attempt` + 1, in milliseconds: what
  -- the API asked for, if it did; otherwise 2, 4, 8, ... seconds.
  retry_delay(response: Http_Response, attempt: Integer): Integer do
    result := retry_delay_for(attempt)
    across response.headers().keys() as key do
      if key.to_lower() = "retry-after" then
        do
          result := (response.headers().get(key).trim().to_integer * 1000).max(1000)
        rescue
          result := retry_delay_for(attempt)
        end
      end
    end
  end

  retry_delay_for(attempt: Integer): Integer do
    result := 1000
    repeat attempt.min(5) do
      result := result * 2
    end
  end

  -- An API error as a sentence: "Claude API error 401
  -- (authentication_error): invalid x-api-key".
  api_error(response: Http_Response): String do
    result := "Claude API error " + response.status().to_string
    do
      let body: Json_Value := json_value(response.body())
      result := result + " (" + body.get("error").get("type").as_string + "): "
        + body.get("error").get("message").as_string
    rescue
      result := result + ": " + preview(response.body(), 200)
    end
  end
end

-- Whether the environment holds credentials for the Claude API.
function has_credentials(): Boolean do
  result := (program.has_env("ANTHROPIC_API_KEY") and program.env("ANTHROPIC_API_KEY") /= "")
    or (program.has_env("ANTHROPIC_AUTH_TOKEN") and program.env("ANTHROPIC_AUTH_TOKEN") /= "")
end
\end{continued}

Header names are not case-sensitive in HTTP, so \code{retry_delay} compares
them in lower case.

\section{A pretend Claude}

We can try all of this without a key and without the network: Nex can be
an HTTP server too. This try-it starts a pretend Messages API on this
computer, answering like Claude would. The first request gets a 529
(``overloaded''), to see the retry. After that it reads the conversation it
is sent: with one message, it asks to list the files; with more, it says
it is done. Then we point a \code{Claude_Model} at it, and give it to the
agent:

\begin{tryit}{try_pretend_api.nex}
intern Agent
intern Claude_Model
intern net/Http_Server

let requests: Integer := 0
let server: Http_Server := create Http_Server.make(0)
server.post("/v1/messages", fn (request: Http_Request): Http_Server_Response do
  requests := requests + 1
  let body: Json_Value := json_value(request.body())
  if requests = 1 then
    result := create Http_Server_Response.make(529,
      '{"type": "error", "error": {"type": "overloaded_error", "message": "Overloaded"}}',
      {"retry-after": "1"})
  elseif body.get("messages").length = 1 then
    result := create Http_Server_Response.text('{"content": [{"type": "text", "text": "I will look around."}, {"type": "tool_use", "id": "toolu_1", "name": "list_files", "input": {}}], "stop_reason": "tool_use", "usage": {"input_tokens": 2150, "output_tokens": 40}}')
  else
    result := create Http_Server_Response.text('{"content": [{"type": "text", "text": "There are '
      + body.get("tools").length.to_string + ' tools, and the model is '
      + body.get("model").as_string + '. Done."}], "stop_reason": "end_turn", "usage": {"input_tokens": 12, "cache_read_input_tokens": 2190, "output_tokens": 20}}')
  end
end)
server.start()

let ws: Workspace := fresh_workspace(user_path("sample_project"), user_path("scratch"))
let claude: Claude_Model := create Claude_Model.make("not-a-real-key", false, "claude-opus-5-5", "high")
claude.set_base_url("http://127.0.0.1:" + server.port.to_string)
console.print_line(claude.headers().to_string)
let agent: Agent := create Agent.make(claude, standard_tools(ws), create Allow_Everything.make(),
  create Reporter.make(false, false))
agent.run("Look around.")
server.stop()
\end{tryit}
\tryoutput{try_pretend_api}

That was the agent talking HTTP to something that behaves like Claude: a
busy API, a tool call, a final answer, and a token count with the second
request read mostly from the cache. The server's handler is a \code{fn},
a function written in place, and it can use and change the variables
around it, such as \code{requests}.

\section{With a key}

To use Claude itself, get an API key from \code{console.anthropic.com} and
put it in the environment:

\begin{shell}
export ANTHROPIC_API_KEY=sk-ant-...
\end{shell}

Then \code{create Claude_Model.from_environment("claude-opus-5-5",
"high")} is a model like any other. The program in
Chapter~\ref{ch:program} does exactly that. First, though, a way to keep
what Claude did.
